\section*{Hoofdstuk \ref{ch:g10:lines} --- Vergelijkingen van rechten en lineaire stelsels}

\begin{solution}{exo:g10:lines:1}
$y = 2x - 3$: \omterm{def:g10:reffunc:affine}{richtingscoëfficiënt} $2$, $y$-afsnede $-3$.
$y = -\frac13 x + 1$: \omterm{def:g10:reffunc:affine}{richtingscoëfficiënt} $-\frac13$, $y$-afsnede $1$.
$y = 4$: \omterm{def:g10:reffunc:affine}{richtingscoëfficiënt} $0$, $y$-afsnede $4$ (horizontale rechte).
$x = -2$: verticale rechte, geen \omterm{def:g10:reffunc:affine}{richtingscoëfficiënt} en geen \omterm{thm:g10:lines:reduced}{herleide
vergelijking}.
\end{solution}

\begin{solution}{exo:g10:lines:2}
$3 \times 2 + 1 = 7$: ja, $A$ ligt op de rechte.
$3 \times (-1) + 1 = -2 \neq -1$: $B$ niet.
$3 \times 4 + 1 = 13$: ja, $C$ ligt op de rechte.
\end{solution}

\begin{solution}{exo:g10:lines:3}
(a) $m = \frac{8 - 2}{3 - 0} = 2$ en $p = y_A = 2$ (het punt $A$ ligt op de
$y$-as): $y = 2x + 2$.

(b) $m = \frac{-4 - 5}{2 - (-1)} = \frac{-9}{3} = -3$; daarna geeft
$5 = -3 \times (-1) + p$ dat $p = 2$: $y = -3x + 2$.

(c) $x_A = x_B = 4$: verticale rechte $x = 4$.
\end{solution}

\begin{solution}{exo:g10:lines:4}
$y = \frac{6x+2}{2} = 3x + 1$. De rechten met \omterm{def:g10:reffunc:affine}{richtingscoëfficiënt} $3$ zijn
$y = 3x - 1$, $y = 3x + 5$ en $y = 3x + 1$: die drie zijn onderling
evenwijdig; geen twee ervan vallen samen (verschillende $y$-afsneden), en
geen enkele is evenwijdig met $y = -3x + 2$.
\end{solution}

\begin{solution}{exo:g10:lines:5}
Eerste stelsel: vul $y = 2x - 1$ in $3x + y = 9$ in: $3x + 2x - 1 = 9$, dus
$5x = 10$, $x = 2$, en dan $y = 3$. Oplossing: $(2, 3)$.

Tweede stelsel: uit $x - y = 4$ volgt $x = y + 4$. Dan is
$2(y + 4) + 3y = 3$, dus $5y = -5$, $y = -1$, en dan $x = 3$. Oplossing:
$(3, -1)$.
\end{solution}

\begin{solution}{exo:g10:lines:6}
Eerste stelsel: de \omterm{def:g10:algebra:equation}{vergelijkingen} optellen werkt $y$ weg: $2x = 14$, $x = 7$;
daarna geeft $7 + y = 10$ dat $y = 3$. Oplossing: $(7, 3)$.

Tweede stelsel: vermenigvuldig de eerste \omterm{def:g10:algebra:equation}{vergelijking} met $5$ en de tweede
met $-2$:
\[
\begin{cases}
15x + 10y = 60 \\
-4x - 10y = -16
\end{cases}
\]
Optellen: $11x = 44$, dus $x = 4$; daarna geeft $3 \times 4 + 2y = 12$ dat
$y = 0$. Oplossing: $(4, 0)$.
\end{solution}

\begin{solution}{exo:g10:lines:7}
Bereken telkens $ab' - a'b$.

Eerste: $2 \times 2 - (-4) \times (-1) = 4 - 4 = 0$, en de tweede
\omterm{def:g10:algebra:equation}{vergelijking} is $-2$ keer de eerste: twee keer dezelfde rechte, oneindig veel
oplossingen.

Tweede: opnieuw $ab' - a'b = 0$, maar de rechterleden staan niet in de
verhouding $-2$ ($1 \neq -6$): twee verschillende evenwijdige rechten, geen
oplossing.

Derde: $2 \times 1 - 1 \times (-1) = 3 \neq 0$: precies één oplossing.
\end{solution}

\begin{solution}{exo:g10:lines:8}
Zij $x$ het aantal kinderkaartjes en $y$ het aantal volwassenenkaartjes:
\[
\begin{cases}
x + y = 200 \\
5x + 9y = 1352 .
\end{cases}
\]
Uit de eerste \omterm{def:g10:algebra:equation}{vergelijking} volgt $x = 200 - y$; invullen geeft
$5(200 - y) + 9y = 1352$, dus $1000 + 4y = 1352$, $y = 88$, en dan
$x = 112$. Dus $112$ kinderkaartjes en $88$ volwassenenkaartjes. Controle:
$5 \times 112 + 9 \times 88 = 560 + 792 = 1352$.
\end{solution}

\begin{solution}{exo:g10:lines:9}
\emph{1.} Evenwijdig betekent dezelfde \omterm{def:g10:reffunc:affine}{richtingscoëfficiënt} $2$:
$y = 2x + p$ met $4 = 2 \times 1 + p$, dus $p = 2$: $d'\colon y = 2x + 2$.

\emph{2.} Op de $x$-as is $y = 0$: $2x + 2 = 0$ geeft $x = -1$. Het snijpunt
is $(-1, 0)$.
\end{solution}

\begin{solution}{exo:g10:lines:10}
\emph{1.} \omterm{prop:g10:coordgeom:midpoint}{Midden} van $[BC]$: $A' = \left(\frac{6+2}{2},
\frac{2+4}{2}\right) = (4, 3)$. De zwaartelijn uit $A(0,0)$ heeft
\omterm{def:g10:reffunc:affine}{richtingscoëfficiënt} $\frac{3 - 0}{4 - 0} = \frac34$: \omterm{def:g10:algebra:equation}{vergelijking}
$y = \frac34 x$.

\omterm{prop:g10:coordgeom:midpoint}{Midden} van $[AC]$: $B' = (1, 2)$. De zwaartelijn uit $B(6, 2)$ heeft
\omterm{def:g10:reffunc:affine}{richtingscoëfficiënt} $\frac{2 - 2}{1 - 6} = 0$: het is de horizontale rechte
$y = 2$.

\emph{2.} Snijpunt: $\frac34 x = 2$ geeft $x = \frac83$, dus
$G\left(\frac83, 2\right)$. De derde zwaartelijn verbindt $C(2,4)$ met het
\omterm{prop:g10:coordgeom:midpoint}{midden} van $[AB]$, $C' = (3, 1)$; haar \omterm{def:g10:reffunc:affine}{richtingscoëfficiënt} is
$\frac{1 - 4}{3 - 2} = -3$, \omterm{def:g10:algebra:equation}{vergelijking} $y = -3x + 10$. In $x = \frac83$:
$y = -8 + 10 = 2$: $G$ ligt erop. En inderdaad is
$G = \left(\frac{0 + 6 + 2}{3}, \frac{0 + 2 + 4}{3}\right)$: het zwaartepunt
middelt de \omterm{def:g10:coordgeom:system}{coördinaten} van de hoekpunten.
\end{solution}

\begin{solution}{pb:g10:lines:1}
\textbf{1.} \omterm{def:g10:reffunc:affine}{Richtingscoëfficiënt} $\frac{8 - 2}{3 - 1} = 3$; door $(1, 2)$:
$y = 3x - 1$.

\textbf{2.} $y = 3x - 1$ en $y = 3x + 4$: dezelfde \omterm{def:g10:reffunc:affine}{richtingscoëfficiënt}, dus
evenwijdig (en verschillend). $y = 3x - 1$ snijden met $y = -x + 7$:
$3x - 1 = -x + 7$ geeft $x = 2$, $y = 5$: het punt $(2, 5)$.

\textbf{3.} Beide punten delen $x = 2$: de rechte is verticaal, met
\omterm{def:g10:algebra:equation}{vergelijking} $x = 2$. Een herleide vorm $y = mx + p$ beantwoordt de vraag
``één $y$ per $x$'', en dat weigert een verticale rechte: haar
\omterm{def:g10:reffunc:affine}{richtingscoëfficiënt} zou oneindig zijn.

\textbf{4.} $3 \times 4 - 1 = 11$: ja, $(4, 11)$ ligt op de rechte.
$3 \times (-1) - 1 = -4 \neq -5$: nee.

\textbf{5.} $y = -2x + 3$; snijdt de assen in $(0, 3)$ en
$\left(\frac32, 0\right)$. Oppervlakte van de driehoek:
$\frac12 \times \frac32 \times 3 = \frac94 = 2.25$.

\textbf{6.} Invullen: $3x + 2x - 3 = 7$, dus $x = 2$ en $y = 1$: oplossing
$(2, 1)$.

\textbf{7.} Optellen: $7x = 7$, dus $x = 1$; daarna $3y = 6$, $y = 2$:
oplossing $(1, 2)$.

\textbf{8.} Eerste stelsel: de tweede \omterm{def:g10:algebra:equation}{vergelijking} is het dubbele van de
eerste: één rechte die twee keer geteld wordt --- oneindig veel oplossingen
(alle punten van $2x - y = 3$). Tweede stelsel: dezelfde evenredige
linkerleden, maar niet-evenredige rechterleden ($2 \times 3 = 6 \neq 1$):
twee evenwijdige rechten --- geen oplossing. Driedeling: verschillende
\omterm{def:g10:reffunc:affine}{richtingscoëfficiënten} $\to$ één snijpunt; gelijke \omterm{def:g10:reffunc:affine}{richtingscoëfficiënten} met
verschillende $y$-afsneden $\to$ evenwijdig, geen oplossing; alles gelijk
$\to$ dezelfde rechte, oneindig veel.

\textbf{9.} De richtingsvectoren van de rechten zijn $(-b, a)$ en $(-d, c)$
(of: de \omterm{def:g10:reffunc:affine}{richtingscoëfficiënten} $-\frac ab$ en $-\frac cd$), en ze zijn
\omterm{def:g10:vectors:collinear}{collineair} precies wanneer $ad - bc = 0$ --- de determinant van
\cref{pb:g10:vectors:1} in een nieuwe betrekking. Waarden: vraag 7:
$2 \times (-3) - 3 \times 5 = -21 \neq 0$: één oplossing. Vraag 8:
$2 \times (-2) - (-1) \times 4 = 0$ voor beide: evenwijdig of samenvallend,
en de constanten scheiden de twee gevallen.

\textbf{10.} $ad - bc = 6 - 2m$: nul voor $m = 3$. Daar staat $x + 3y = 3$
tegenover $2x + 6y = 7$: het dubbele van de eerste geeft
$2x + 6y = 6 \neq 7$: evenwijdig, dus \emph{geen oplossing}. Voor
$m \neq 3$ levert eliminatie de enige oplossing $y = \frac{1}{6 - 2m}$, en
daarna $x = 3 - \frac{m}{6 - 2m}$.

\textbf{11.} $f + k = 35$, $2f + 4k = 94$: halveren geeft $f + 2k = 47$, dus
$k = 12$ en $f = 23$: drieëntwintig fazanten en twaalf konijnen --- het
antwoord van de \emph{Negen Hoofdstukken} zelf.

\textbf{12.} $x + y = 30$ en $0.6x + 0.2y = 0.35 \times 30 = 10.5$; $0.2$
keer de eerste ervan aftrekken geeft $0.4x = 4.5$: $x = 11.25$ L siroop en
$y = 18.75$ L drank.

\textbf{13.} Cijfers $a$, $b$: $a + b = 12$ en
$(10a + b) - (10b + a) = 9(a - b) = 18$, dus $a - b = 2$: $a = 7$, $b = 5$:
het getal is $75$.

\textbf{14.} Optellen: $5k + 5c = 18$, dus $k + c = 3.60$ --- één koffie plus
één croissant. Aftrekken: $k - c = 0.20$. Bijgevolg is $k = 1.90$ en
$c = 1.70$ euro: de symmetrische kortere weg loste het op zonder één
substitutie.

\textbf{15.} Het dubbele van de eerste bewering geeft $4$ appels $+ 2$ peren
$= 10$ euro, terwijl de kraam $9$ vraagt: het stelsel is strijdig ---
evenwijdige rechten, lege oplossingsverzameling. Oordeel: de twee
prijsbeweringen kunnen niet allebei kloppen; ofwel zit er een korting
verborgen, ofwel de rekenfout.

\textbf{16.} In de oorsprong: $0 > 3 \times 0 - 1 = -1$: waar, dus ligt
$(0,0)$ in het gebied $y > 3x - 1$ (boven de rechte). Toets een willekeurig
punt door het in te vullen: boven of onder, naargelang de ongelijkheid waar
is of niet.

\textbf{17.} Hoekpunten: $(0, 0)$; $(4, 0)$ (assen en $x + y = 4$);
$(0, 1)$ (as en $y = 2x + 1$); en $x + y = 4$ met $y = 2x + 1$: $x = 1$,
$y = 3$: $(1, 3)$.

\textbf{18.} $P(0,0) = 0$; $P(4, 0) = 12$; $P(0, 1) = 2$; $P(1, 3) = 9$.
\omterm{def:g10:functions:extrema}{Maximum} in het hoekpunt $(4, 0)$: het plan $x = 4$, $y = 0$ levert $12$ op
--- de hoekpunten beslissen, zoals de strijkende evenwijdige rechten
$3x + 2y = c$ zichtbaar maken.

\textbf{19.} $1.5x + 3 = 2x$ geeft $x = 6$, $y = 12$: bij zes kilometer
rekenen beide taxi's twaalf euro --- het omslagpunt uit de weekendopgave over
de twee thermometers in het onderbouwvolume, nu de enige oplossing van een
stelsel.

\textbf{20.} De oplossingen van een lineaire \omterm{def:g10:algebra:equation}{vergelijking} tekenen een rechte;
een stelsel legt twee rechten over elkaar, en $ad - bc$ zegt in één oogopslag
of ze elkaar één keer snijden (niet nul) dan wel in de gevallen
evenwijdig-of-samenvallend belanden (nul). Stapel in plaats daarvan
ongelijkheden op, en de oplossingen vormen een veelhoekig gebied waarvan de
hoekpunten elk lineair optimum dragen. De lineaire algebra veralgemeent de
determinant naar willekeurig veel onbekenden; het lineair programmeren maakt
industrie van de hoekpunten.
\end{solution}
